% =====================================================================
% 6  Conclusion                        Folien 15-16 | 1:45 min | bis 19:45
% Inhalt: Bericht Kap. 9 (Aufgabe 6: is there reliable evidence on the
% true effect size?). Restliche ca. 15 s = Puffer.
% =====================================================================
\section{Conclusion}

% ---------------------------------------------------------------------
\begin{frame}{Conclusion: is the evidence reliable?}
  \timing{1:30}{19:30}
  \begin{itemize}
    \item \textbf{Existence:} robust -- all pooled and corrected estimates
      are below 1, even the most conservative (PET, combined:
      \maORpetComb{} \maCIpetComb)
    \item \textbf{Magnitude:} uncertain -- estimates range from \maORsub{}
      to \maORpetComb; $I^2 = \maIsqRe\,\%$, prediction interval \maPIre{}
      includes 1
    \item \textbf{Publication bias:} little sign in the original sample; in
      the combined sample driven by one study and partly by region
    \item \textbf{Power:} high by economics standards (median \maMedPower),
      but likely an upper bound
    \item \textbf{Limitations:} reconstructed counts for three new studies;
      SEs ignore clustering by employer; only the first stage of hiring
  \end{itemize}
  \keymsg{Ethnic discrimination in hiring is robust, but its size varies
    strongly across countries and groups -- a single true effect is hard to
    justify.}
\end{frame}
% Sprechnotizen (Folie 15):
% - Pro Zeile ein Satz, sonst reichen 1:30 nicht.
% - Magnitude: Spannweite von 0.55 (RE, Subgruppen) bis 0.87 (PET, kombiniert).
%   Das Prediction Interval schliesst 1 ein -> in manchen Arbeitsmaerkten
%   kaum oder keine Diskriminierung.
% - Limitations: Bei Zeitdruck nur zwei nennen.

% ---------------------------------------------------------------------
\begin{frame}[c]
  \timing{0:15}{19:45}
  \centering
  {\usebeamercolor[fg]{structure}\Large\bfseries Thank you!}\\[1em]
  Questions and discussion\\[2.5em]
  {\small Kevin de Silva \quad\textbar\quad Vanessa Allenspach}
\end{frame}
% Vor dem Mailen an den Prof: in slides.tex \showtimingfalse setzen.
